% ---
% id: transmitter/def:source-spec
% kind: def
% title: "Source specification"
% summary: >-
%   Alice's alphabet, retained system with sectors, measurement and signal
%   ensemble are declared, with binary64 numbers and exact decisions, or
%   defined from declared data or a source model. Trusted-source mode adds a
%   marginal and, above dimension one, a certificate; untrusted-source mode
%   has neither.
% prominence: standard
% status: implemented
% uses: []
% mentions: [def:source-segment, def:srcert, rem:source-mode, lem:matched-pinching, ass:ens, lem:ensemble-presentation, rem:certificate-forms, rem:da1]
% assumes: []
% refs: []
% provenance: newly-derived
% implements: certiqs_key.transmitter.artifact:SourceSpecification
% lean:
% notes: "This definition and def:srcert define each other (trusted-source mode is defined by the certificate, the certificate by this definition's data), but the graph carries only the edge def:srcert -> def:source-spec: a change to def:srcert's conditions reaches this definition by hand. Also realised by certiqs_key.transmitter.artifact:emit_artifact, which decides the conditions."
% ---

Alice's side of a protocol is described by what she prepares and what she
keeps. The definition lists the data a specification declares, the
conditions on them, and the objects that are defined from declared data
rather than declared themselves; its closing number model says how declared
data are given and how every condition is decided. In trusted-source mode the specification asserts that Alice's retained system
$A$ never reaches the adversary, and a \tref{source-replacement certificate}
(Definition~\ref{def:srcert}) describes her preparation as her measurement
$\{M_x\}$ of $A$. In untrusted-source mode no certificate ties $A$ or
$\{M_x\}$ to the signal ensemble: $\{M_x\}$ is Alice's measurement of a system
the adversary may have prepared, and the mode is the conservative
relaxation of trusted-source mode (Remark~\ref{rem:source-mode}). The sector
structure on $A$ is what lets a certificate be declared in \tref{block form}
(Definition~\ref{def:srcert}), which a source that first draws a sector and
then prepares a state in it needs (Remark~\ref{rem:certificate-forms}).
At $d_A=1$ the two modes coincide (Remark~\ref{rem:da1}).

\begin{definition}\label{def:source-spec}
A \term[The declared description of Alice's side of a protocol: alphabet, retained system with sectors, measurement, signal ensemble and source mode.]{source specification} $\Spec_{\mathrm{src}}$ consists of the
following data. A finite non-empty alphabet $\Xset$ with a declared order. A
finite-dimensional system $A$ with a declared computational basis,
$d_A\coloneqq\dim\Hs{A}$, and a declared \term[A decomposition of Alice's retained system into orthogonal sectors, each spanned by declared basis vectors.]{sector structure} on it: an
ordered finite index set $\mathbb{T}$ and a decomposition
$\Hs{A}=\bigoplus_{t\in\mathbb{T}}\Hs{A_t}$ into subspaces each spanned by
vectors of the declared basis, $d_{A_t}\coloneqq\dim\Hs{A_t}\ge1$, with
projectors $\Pi^{A}_t$; the \term[The default sector structure, whose single sector is the whole retained system.]{one-sector structure}, $\mathbb{T}$ a
single index and $A_t=A$, is the default. A POVM $\{M_x\}_{x\in\Xset}$ on
$A$ presented as a block list against that structure,
$M_x=\bigoplus_{t}M^{(t)}_x$ with $M^{(t)}_x\in\Lin{A_t}$, each block
Hermitian and positive semidefinite, and
$\sum_{x\in\Xset}M^{(t)}_x=\idon{A_t}$ for every $t$ with the sum taken in
the declared order. The \term[The weights and signal states, one per letter, that Alice's source is specified to emit.]{specified signal ensemble}
$\{p_x,\rho_x\}_{x\in\Xset}$ on the \term[The system Alice sends into the channel, carrying the specified signal ensemble.]{transmitted system} $\Asig$, with
$p_x\ge0$, $\sum_xp_x=1$ and $\rho_x$ a state on $\Asig$. With one sector
the block list $M_x$ is the matrix $M_x$, and every condition on its blocks
is the same condition on that matrix.

\emph{How the specified ensemble is declared.} In one of two ways.
(a) \emph{As matrices}, for a finite-dimensional $\Asig$: by its
\term[The positive operators, weight times state, one per letter, through which an ensemble is declared and read.]{weighted members} $\sigma_x\in\Pos{\Asig}$, with
$p_x\coloneqq\tr\sigma_x$, $\sum_xp_x=1$ in the declared order, and
$\rho_x\coloneqq\sigma_x/p_x$ where $p_x>0$. A declaration by the pair
$(p_x,\rho_x)$, $\rho_x\in\Seq{\Asig}$, is the declaration of the weighted
member $\sigma_x=p_x\rho_x$ with the product formed exactly, and it requires
$\rho_x\in\Seq{\Asig}$ of \emph{every} letter, a letter with $p_x=0$
included. Every condition of this module that reads an admitted ensemble
reads it through its weighted members only.
(b) \emph{By a} \term[A submodule's named definition of the signal ensemble as mathematical objects, possibly infinite dimensional, rather than as matrices.]{source model}, for a transmitted system $\Asig$ whose
Hilbert space is separable, possibly infinite dimensional: by a named
definition of a submodule of this module that fixes the weights $p_x$ and the
states $\rho_x$ as mathematical objects and not as matrices, together with
the physical-source assumption which that submodule declares in the place of
Assumption~\ref{ass:ens}. A specification declared by a source model is
presented with a non-empty \tref{source segment}
(Definition~\ref{def:source-segment}): what a consumer receives is never the
model but the \tref{virtual ensemble}, which is finite dimensional.

\emph{The two modes.} The specification is declared in one of two modes. In
\term[The source mode in which Alice's retained system never reaches the adversary, with a declared marginal and, above dimension one, a certificate.]{trusted-source mode} it additionally carries a marginal
$\tau_A\in\Seq{A}$, presented as a block list
$\tau_A=\bigoplus_t\tau^{(t)}_A$ with $\tau^{(t)}_A\in\Pos{A_t}$ and
$\sum_t\tr\tau^{(t)}_A=1$, and, at $d_A\ge2$, a source-replacement
certificate in the sense of Definition~\ref{def:srcert}, written against the
\emph{virtual} ensemble $\{p_x,\rvirt_x\}$ on the \tref{virtual transmitted system}
$\Avirt$: the finite-dimensional ensemble the source segment starts from, whose
channels produce the specified ensemble from it
(Definition~\ref{def:source-segment}\,(S1), (S3)); when the source segment is
empty, $\Avirt=\Asig$ and $\rvirt_x=\rho_x$. The certificate takes one of the
three forms of Definition~\ref{def:srcert}: a rank-one projector (the \tref{pure
form}), positive sector blocks (the block form), or an \tref{ensemble presentation},
which declares only the weighted virtual ensemble, in blocks along a matched
sector structure of $\Avirt$, and defines the certificate from it. A sector structure with more than one sector is declared only in
trusted-source mode and only together with a certificate in block form or an
ensemble presentation. \emph{With an ensemble presentation} the specification
declares neither $A$, its basis and its sector structure, nor $\{M_x\}$, nor
$\tau_A$: they are constituents of the specification that \eqref{eq:enspres}
defines from the declared virtual ensemble, the conditions this definition
places on them hold by Lemma~\ref{lem:ensemble-presentation}\,(i), and of
them only $\tau_A$ is formed as numbers, that is, computed entry by entry from
the declared blocks (Lemma~\ref{lem:ensemble-presentation}\,(iii));
$\{M_x\}$ and the certificate are defined and never computed. At $d_A=1$ no
certificate is demanded and $\tau_A=[1]$ is the only marginal, so that there
the two modes impose the same condition. In
\term[The source mode with no certificate and no marginal, in which the adversary may have prepared Alice's retained system.]{untrusted-source mode} the specification carries no certificate and no
marginal, written
$\tau_A=\nomarg$; the specification with
$d_A=1$ and a one-letter $\Xset$ is its degenerate case, in which Alice
prepares nothing that this module describes.

\emph{What the sector structure is not.} It is a declaration about
$\{M_x\}$ and $\tau_A$, which are block diagonal by their presentation. It
is no statement about any operator on $A$ together with a further system,
and no consumer may read it as one: that a joint state is block diagonal
along $\mathbb{T}$ is either proved of the states a certificate reaches
(Lemma~\ref{lem:matched-pinching}) or licensed by a theorem of the consuming
module that restricts its own feasible set, under that module's own
preconditions, and is never inferred from the presentation of $\tau_A$ or
$\{M_x\}$.

\emph{The number model.} Every declared entry is a binary64 number, hence a
dyadic rational, and is admissible exactly as declared. A source model
(case (b)) declares no entries: its weights and states are mathematical
objects, and what is decided of such a specification is decided on the
entries it does declare, those of its virtual ensemble among them. Every
condition above is decided \emph{exactly} on the declared entries: a
quantity that a condition compares --- a sum in the declared order, a
product, a partial
trace --- is \emph{formed} in exact arithmetic on the declared entries as
well as compared exactly, so that the identity decided is the identity
written and rounding enters nowhere. That the operands are binary64 is not
what makes a decision exact; a product of two binary64 numbers need not be
one. Positive semidefiniteness is decided by an exact decision on the
declared entries and never by a floating-point eigenvalue.
\end{definition}
